\documentclass[11pt]{article}
\usepackage[margin=1in]{geometry}
\usepackage{times}
\usepackage{amsmath,amssymb}
\usepackage{graphicx}
\usepackage{booktabs}
\usepackage{multirow}
\usepackage{caption}
\usepackage{enumitem}
\usepackage{longtable}
\usepackage{array}
\usepackage{url}
\usepackage[hidelinks]{hyperref}

\title{Supplementary Material for\\Schema-Aware Grounding Effects in PostGIS\\
Natural-Language-to-SQL: A Subset-Decomposed Evaluation\\Across Eleven LLMs}
\author{}
\date{}

\begin{document}
\maketitle

\noindent This supplement collects two pieces of material referenced from the main text but moved out of the main body to keep the manuscript focused. Cross-references from the main manuscript appear as ``Supplement~S1/S2''.

\begin{itemize}[leftmargin=*]
\item \textbf{S1.} Full baseline framework evaluation: BIRD warehouse track design-set and held-out, cross-lingual per-question re-audit, DIN-SQL external comparison, full reproducibility inventory, and the cross-model-family $2{\times}2$ factorial that motivated the broader cross-family panel reported in main-text Section~6.
\item \textbf{S2.} End-to-end worked example tracing one Chongqing benchmark question through every pipeline stage (referenced from main text \S3).
\end{itemize}

\section*{S1. Baseline framework evaluation (full version)}

The main text \S5 provides a compact summary of framework behaviour on the GIS Spatial set, BIRD warehouse track, DIN-SQL comparison, and cross-lingual probe; this section provides the full multi-section version with all per-question paired McNemar tables. These framework-level runs are supporting context and use their original protocols. The primary cross-family grounding-effect analysis in main-text Section~6 uses question-level majority vote followed by a single exact two-sided McNemar test on the per-question table.

\subsection*{S1.1 BIRD mini\_dev design-set and held-out details}

This section expands the warehouse-track stub in the main text \S5.
main text. We report the BIRD evaluation in two parts: a \emph{design set} of 108 paired
questions on which grounding rules were derived by error attribution, and a disjoint
\emph{held-out set} of 150 paired questions on which the same rules are tested without
further tuning.

\paragraph{Round-1 grounding (design set).} Round-1 adds three rule sections to the
prompt: (i)~multi-hop join-path discovery bridging fact-to-fact tables through pivot
dimensions; (ii)~aggregation semantics clarifying \texttt{COUNT(*)} vs.\ \texttt{COUNT(col)}
vs.\ \texttt{COUNT(DISTINCT col)}; (iii)~temporal handling for TEXT-typed date columns
(prefix comparison, \texttt{SUBSTR} extraction). On 101 paired questions it gave
EX$=0.564$ vs.\ baseline $0.515$ ($+0.050$, paired McNemar two-sided $p=0.2266$,
$b=3$, $c=8$).

\paragraph{Round-2 grounding (design set).} Error attribution on round-1 failures
identified three patterns: 16 of 46 both-fails differed from the gold SQL only in
\texttt{DISTINCT} placement (SELECT list dimensions duplicated through one-to-many
JOINs); 4 regressions were caused by over-JOINing when a single table sufficed;
multiple \texttt{student\_club} questions used \texttt{||} concatenation where the
gold returned first/last name as separate columns. Round-2 adds three targeted rules:
DISTINCT-after-one-to-many-JOIN, avoid-over-JOIN, and output-column format (full name
as separate columns; \texttt{LIMIT 1} via \texttt{ORDER BY} rather than subquery with
\texttt{MAX}/\texttt{MIN}).

Table~\ref{tab:bird-supp} reports the round-2 result on 108 paired design questions and
on the 150-question held-out set. Significance on the design set alone does not
establish generalisation because the rules were derived from the same 108 items.

\begin{table}[h]
\centering\footnotesize
\caption{BIRD warehouse results across rounds and splits.}
\label{tab:bird-supp}
{\setlength{\tabcolsep}{3pt}
\begin{tabular}{lc>{\raggedright\arraybackslash}p{2.8cm}>{\raggedright\arraybackslash}p{2.8cm}>{\raggedright\arraybackslash}p{3.0cm}}
\toprule
Set & $N$ & Baseline EX [95\% CI] & Full EX [95\% CI] & $\Delta$ [95\% CI] \\
\midrule
Round-1 grounding (design, subset) & 101 & 0.515 & 0.564 & $+$0.050 ($p=0.227$, NS) \\
\textbf{Round-2 design set (tuning)} & \textbf{108} & \textbf{0.500 [0.41, 0.59]} & \textbf{0.593 [0.50, 0.68]} & $+$\textbf{0.093} [0.02, 0.16] ($p=0.021$) \\
\textbf{Round-2 held-out (primary)}  & \textbf{150} & \textbf{0.473 [0.40, 0.55]} & \textbf{0.507 [0.43, 0.59]} & $+$\textbf{0.033} [$-$0.03, 0.09] ($p=0.383$, NS) \\
\midrule
Validity (R2 design) & 108 & 0.981 & 0.991 & $+$0.010 \\
Validity (R2 held-out) & 150 & 0.980 & 1.000 & $+$0.020 \\
Mean tokens (R2 design) & 108 & 1{,}010 & 7{,}975 & $7.9\times$ \\
\bottomrule
\end{tabular}
}
\end{table}

\paragraph{Interpretation.} The round-2 design-set gain of $+0.093$ with $p=0.0213$
should be read as a \emph{development-set tuning effect}: the rules were derived by
error attribution on this same 108-item sample, so significance on this set documents
that the rules address the specific failure modes observed there. The primary
warehouse-side significance test is therefore the held-out evaluation, where the
same rules yield only $+0.033$ with 95\% paired-proportion CI $[-0.03, 0.09]$ and
$p=0.3833$ on 21 discordant pairs ($b=8$, $c=13$); direction is preserved, magnitude
drops by about two-thirds, and the CI comfortably covers zero. Two readings are
consistent with this pattern: (a)~part of the design-set gain is in-sample
specialisation to the 108-item mix; (b)~at $n=150$ the effect-size-to-variance ratio
remains low even if a small true effect exists. We therefore do not claim a
BIRD-wide improvement from round-2 grounding. Validity improves on both sets
($0.981\to 0.991$ design, $0.980\to 1.000$ held-out), which we attribute to the
postprocessor's safety checks rather than to rule content. Token cost on the design
set is $7.9\times$ baseline, down from $\sim$32$\times$ in the earlier agent-loop
configuration.

Representative design-set wins: Q1155 (``\emph{List patient ID, sex and birthday with
LDH beyond normal range}'') --- round-1 omitted \texttt{DISTINCT}, round-2 added it;
Q1334 (``\emph{full names of members from Illinois}'') --- round-1 used string
concatenation, round-2 returned separate columns; Q1506 --- round-1 missed
\texttt{DISTINCT}, round-2 added it.

\paragraph{Relation to the earlier multi-pass evaluation.} A multi-pass agent-loop
run on a 500-question BIRD subset (without round-2 rules) previously reported
baseline $0.474$, full $0.501$, $\Delta=+0.027$ with $p=0.136$. We retain this only
as context for the token-cost comparison; the 150-question held-out evaluation is
the appropriate independent test of the single-pass grounded pipeline.

\paragraph{DIN-SQL older 495q comparison and held-out paired run.} For the BIRD
track, DIN-SQL was previously evaluated on a 495-question single-pass run as
context for token cost; on that older sample, all three methods (Baseline $0.474$,
DIN-SQL $0.482$, Full $0.501$) cluster between $0.47$ and $0.51$ and the
full-vs-DIN-SQL difference is $+0.019$ ($p=0.382$, NS). The DIN-SQL row in
Table~6 of the main text reports that older comparison only. We additionally
run a paired DIN-SQL comparison on the round-2 150-question held-out set against
the same Full pipeline configuration; the result is reported in Section~S1.4
(Table~\ref{tab:dinsql-bird-heldout}): Full $0.507$ vs.\ DIN-SQL $0.440$, $b=18$,
$c=8$, paired McNemar two-sided exact $p=0.0755$ (marginal). The 495q row is
retained for token-cost continuity; the held-out row is the appropriate independent
DIN-SQL comparison.

\subsection*{S1.2 Cross-lingual re-audit: per-question detail}

Table~\ref{tab:crosslingual-per-question} lists the $50$ BIRD Chinese items, their
reviewer label (ok / fix / drop), the pre-review prediction correctness (0/1), and
the post-review prediction correctness (0/1). This table is machine-generated from
the re-evaluation harness.

% [Rows rendered by scripts/nl2sql_bench_cq/render_crosslingual_supplement.py]
\begin{longtable}{lllccc}
\caption{Per-question reviewer label and pre/post EX on the BIRD Chinese 50-question re-audit.}
\label{tab:crosslingual-per-question} \\
\toprule
qid & db\_id & difficulty & review & pre EX & post EX \\
\midrule
\endfirsthead
\multicolumn{6}{c}{\tablename\ \thetable{} -- continued} \\
\toprule
qid & db\_id & difficulty & review & pre EX & post EX \\
\midrule
\endhead
\midrule
\multicolumn{6}{r}{\emph{continued on next page}} \\
\endfoot
\bottomrule
\endlastfoot
1312 & student\_club & simple & fix & 1 & 1 \\
1317 & student\_club & moderate & fix & 1 & 1 \\
1322 & student\_club & moderate & fix & 0 & 0 \\
1323 & student\_club & moderate & fix & 0 & 0 \\
1331 & student\_club & simple & fix & 0 & 0 \\
1334 & student\_club & simple & ok & 1 & 1 \\
1338 & student\_club & moderate & fix & 0 & 0 \\
1339 & student\_club & challenging & fix & 0 & 0 \\
1340 & student\_club & moderate & fix & 1 & 1 \\
1344 & student\_club & simple & fix & 1 & 1 \\
1346 & student\_club & simple & ok & 1 & 1 \\
1350 & student\_club & moderate & fix & 1 & 1 \\
1351 & student\_club & simple & fix & 1 & 1 \\
1352 & student\_club & moderate & ok & 1 & 1 \\
1356 & student\_club & simple & fix & 1 & 1 \\
1357 & student\_club & simple & ok & 1 & 1 \\
1359 & student\_club & challenging & fix & 1 & 1 \\
1361 & student\_club & simple & ok & 1 & 1 \\
1362 & student\_club & simple & ok & 1 & 1 \\
1368 & student\_club & simple & ok & 1 & 1 \\
1471 & debit\_card\_specializing & simple & ok & 1 & 1 \\
1472 & debit\_card\_specializing & moderate & fix & 1 & 1 \\
1473 & debit\_card\_specializing & moderate & ok & 0 & 0 \\
1476 & debit\_card\_specializing & challenging & fix & 1 & 1 \\
1479 & debit\_card\_specializing & moderate & fix & 0 & 0 \\
1480 & debit\_card\_specializing & moderate & fix & 1 & 1 \\
1481 & debit\_card\_specializing & challenging & fix & 0 & 0 \\
1482 & debit\_card\_specializing & challenging & fix & 0 & 0 \\
1483 & debit\_card\_specializing & simple & ok & 1 & 1 \\
1484 & debit\_card\_specializing & simple & ok & 0 & 0 \\
1486 & debit\_card\_specializing & simple & ok & 0 & 0 \\
1490 & debit\_card\_specializing & moderate & fix & 0 & 0 \\
1493 & debit\_card\_specializing & simple & fix & 1 & 0 \\
1498 & debit\_card\_specializing & simple & fix & 0 & 1 \\
1500 & debit\_card\_specializing & simple & fix & 0 & 0 \\
1501 & debit\_card\_specializing & moderate & ok & 0 & 0 \\
1505 & debit\_card\_specializing & simple & ok & 0 & 0 \\
1506 & debit\_card\_specializing & moderate & fix & 1 & 1 \\
1507 & debit\_card\_specializing & simple & fix & 1 & 1 \\
1509 & debit\_card\_specializing & moderate & ok & 1 & 1 \\
1514 & debit\_card\_specializing & simple & fix & 0 & 0 \\
1515 & debit\_card\_specializing & simple & fix & 1 & 1 \\
1521 & debit\_card\_specializing & moderate & ok & 1 & 1 \\
1524 & debit\_card\_specializing & simple & fix & 1 & 1 \\
1525 & debit\_card\_specializing & simple & ok & 0 & 0 \\
1526 & debit\_card\_specializing & challenging & fix & 0 & 0 \\
1528 & debit\_card\_specializing & simple & fix & 1 & 0 \\
1529 & debit\_card\_specializing & moderate & fix & 0 & 0 \\
1531 & debit\_card\_specializing & moderate & ok & 0 & 0 \\
1533 & debit\_card\_specializing & moderate & fix & 0 & 0 \\
\end{longtable}

\subsection*{S1.3 Cross-model-family full factorial (Gemini vs.\ DeepSeek, $4$ cells $\times$ N$=$3)}

The initial framework evaluation used Gemini~2.5~Flash. To probe whether the observed gain
from semantic grounding and intent-conditioned routing transferred to another LLM
family, we ran a $2{\times}2$ factorial on the full 85-question GIS Spatial
benchmark with model family $\in\{$Gemini~2.5~Flash, DeepSeek-V3$\}$ and
pipeline $\in\{$schema-only baseline, full$\}$. The full pipeline is routed
through a unified model gateway so that the same ADK \texttt{LlmAgent} agent
loop runs on either family without branch-specific code; DeepSeek is reached
via the OpenAI-compatible LiteLLM adapter. Following the main-text Spatial
protocol (Section~3.4 of the manuscript), the full cell is sampled $N{=}3$
times per family and aggregated by majority vote to absorb decoding
stochasticity; each family uses its own provider-default temperature, which
is deployment-realistic. DeepSeek baseline is also run at $N{=}3$ because
DeepSeek's API is not deterministic at temperature~$0$; the Gemini baseline
remains $N{=}1$ because the schema-only \texttt{baseline\_generate} path sets
temperature~$0$ explicitly and Gemini~2.5~Flash is deterministic for that
family. Table~\ref{tab:cross-family-4cell} summarises the four cells;
Table~\ref{tab:cross-family-paired} reports the four paired McNemar tests on
majority-vote predictions.

\begin{table}[h]
\centering
\caption{$4$-cell cross-family factorial on the full 85-question GIS Spatial
benchmark. $\mathrm{EX}$ is execution accuracy over the paired gold-SQL
harness; SD is sample SD across $N=3$ samples within cell. Gemini~full
per-sample EX is reused from the main-text Spatial $N=3$ protocol;
DeepSeek cells are new.}
\label{tab:cross-family-4cell}
\begin{tabular}{lcccc}
\toprule
Cell & EX mean & EX SD & $N$ & Majority-vote EX \\
\midrule
Gemini baseline   & $0.529$ & ---      & $1$ & $0.529$ \\
Gemini full       & $0.663$ & $0.048$  & $3$ & $0.659$ \\
DeepSeek baseline & $0.533$ & $0.007$  & $3$ & $0.529$ \\
DeepSeek full     & $0.459$ & $0.042$  & $3$ & $0.482$ \\
\bottomrule
\end{tabular}
\end{table}

\begin{table}[h]
\centering
\caption{Paired McNemar on majority-vote predictions across the $2{\times}2$
design (two-sided exact binomial). $b$ is the number of qids where the first
cell is correct and the second is incorrect; $c$ is the converse.}
\label{tab:cross-family-paired}
\begin{tabular}{lcccc}
\toprule
Contrast & $b$ & $c$ & $\Delta$EX & Paired $p$ \\
\midrule
Within-family Gemini   (base vs.\ full)     & $8$  & $19$ & $+0.129$ & $0.052$ \\
Within-family DeepSeek (base vs.\ full)     & $18$ & $14$ & $-0.047$ & $0.597$ \\
Cross-family baseline  (Gemini vs.\ DS)     & $0$  & $0$  & $0.000$  & $1.000$ \\
Cross-family full      (Gemini vs.\ DS)     & $22$ & $7$  & $+0.176$ & $0.008$ \\
\bottomrule
\end{tabular}
\end{table}

\paragraph{Reading.} Three findings emerge. \emph{First}, the cross-family
schema-only baseline is exactly matched: both families reach $0.529$ EX with
zero discordant qids ($b=c=0$, $p=1.000$), so the full-pipeline contrast is
not confounded by an a-priori capability gap between the two families on this
benchmark. \emph{Second}, the within-family Gemini grounding gain of $+0.129$
(majority-vote, $p=0.052$) is consistent with the framework-level Spatial result.
\emph{Third}, the same full-pipeline stack --- identical grounding rules,
intent routing, self-correction, and postprocessor --- does \emph{not}
transfer to DeepSeek: the within-family DeepSeek contrast is in the opposite
direction ($-0.047$, $p=0.597$, null), and the cross-family full contrast is
significantly in Gemini's favour ($+0.176$, $p=0.008$). The grounding stack
was authored and error-attributed against Gemini traces, and we interpret the
DeepSeek regression as evidence that semantic grounding plus intent-conditioned
routing, as currently instantiated, carries family-specific prompt
sensitivities. This is an honest negative result for the portability claim:
the mechanism works on the family it was tuned for, does not work on the
unseen family at the same prompt budget, and the next question is whether
family-adapted grounding rules recover the gain.

\paragraph{Protocol details.} Full-pipeline per-sample EX values were:
Gemini $\{0.706, 0.612, 0.671\}$; DeepSeek $\{0.494, 0.412, 0.471\}$.
DeepSeek baseline per-sample EX: $\{0.529, 0.541, 0.529\}$. Temperature
pinning was intentionally \emph{not} applied across families: pinning
temperature~$0$ in a cross-family study would sample a deployment
configuration neither family actually ships; pinning temperature~$1$ would
understate Gemini's deployed determinism. We therefore let each family run
at its provider default and absorbed the decoding variance with $N=3$
sampling per cell, which is the same protocol the main text uses for its
Gemini Spatial headline. All six DeepSeek runs were executed on the same day,
on the same benchmark snapshot, against the same PostGIS instance and the
same gold-SQL harness.

\subsection*{S1.4 DIN-SQL held-out run (150q) and Robustness paired comparison}

DIN-SQL is a published external baseline for text-to-SQL with explicit prompting
stages (schema linking, classification-and-decomposition, self-correction). We
report two paired comparisons of DIN-SQL against the proposed Full pipeline run
with the same external baseline implementation, the same Gemini~2.5~Flash backbone,
and the same execution-comparison harness, against PostgreSQL (BIRD) and PostGIS
(Chongqing) targets.

\paragraph{BIRD 150q held-out.} Table~\ref{tab:dinsql-bird-heldout} pairs DIN-SQL
with the Full pipeline on the 150-question BIRD mini\_dev held-out set --- the same
held-out qid list used by the warehouse evaluation in Section~4.4 of the main text.
The Full pipeline outperforms DIN-SQL by $+0.067$ EX, with discordant counts
$b=18$ in favour of the Full pipeline and $c=8$ in favour of DIN-SQL; paired
McNemar two-sided exact $p=0.0755$ is marginal at $\alpha=0.05$ and consistent
with the small directional effect size on warehouse-style queries reported in
Section~4.4. By BIRD difficulty, DIN-SQL reaches $0.585$ ($24/41$) on simple,
$0.434$ ($33/76$) on moderate, and $0.273$ ($9/33$) on challenging items, so the
DIN-SQL deficit relative to the Full pipeline widens with question difficulty.

\begin{table}[h]
\centering
\caption{Paired DIN-SQL vs.\ Full pipeline on the 150-question BIRD held-out set.
Discordant counts use the convention $b$~=~Full-correct, DIN-incorrect;
$c$~=~DIN-correct, Full-incorrect.}
\label{tab:dinsql-bird-heldout}
\begin{tabular}{lcccc}
\toprule
System & EX & Discordant $b$ & Discordant $c$ & Paired $p$ \\
\midrule
DIN-SQL          & $0.440$ ($66/150$) & ---  & ---  & ---       \\
Full pipeline    & $0.507$ ($76/150$) & $18$ & $8$  & $0.0755$  \\
\bottomrule
\end{tabular}
\end{table}

\paragraph{Robustness 40q.} DIN-SQL was not designed for safety- or refusal-style
queries; this is a stringency test of an external baseline against the Robustness
suite, not a peer benchmark of two safety-aware systems.
Table~\ref{tab:dinsql-cq-robust} pairs DIN-SQL with the Full pipeline on the
40-question Robustness suite. The Full pipeline outperforms DIN-SQL by $+0.700$ EX,
with discordant counts $b=28$ in favour of the Full pipeline and $c=0$ in favour
of DIN-SQL; paired McNemar two-sided exact $p=7.45\times10^{-9}$. Per-category
DIN-SQL accuracy: Security Rejection $3/11$, Anti-Illusion $1/9$, Schema Enforcement
$2/3$, Schema Hallucination $3/8$, Data Tampering Prevention $0/1$, OOM Prevention
$2/8$. The gaps on Security Rejection and Anti-Illusion in particular reflect the
absence of a safety/refusal layer in the DIN-SQL prompt design.

\begin{table}[h]
\centering
\caption{Paired DIN-SQL vs.\ Full pipeline on the 40-question Robustness suite.
Discordant counts use the convention $b$~=~Full-correct, DIN-incorrect;
$c$~=~DIN-correct, Full-incorrect.}
\label{tab:dinsql-cq-robust}
\begin{tabular}{lcccc}
\toprule
System & EX & Discordant $b$ & Discordant $c$ & Paired $p$               \\
\midrule
DIN-SQL          & $0.275$ ($11/40$) & ---  & ---  & ---                   \\
Full pipeline    & $0.975$ ($39/40$) & $28$ & $0$  & $7.45 \times 10^{-9}$ \\
\bottomrule
\end{tabular}
\end{table}

\subsection*{S1.5 Full reproducibility inventory}

\begin{table}[h]
\centering\footnotesize
\caption{Experimental configurations and released artefacts.}
\label{tab:repro-supp}
{\setlength{\tabcolsep}{4pt}
\begin{tabular}{>{\raggedright\arraybackslash}p{3.7cm}>{\raggedright\arraybackslash}p{3.5cm}>{\raggedright\arraybackslash}p{7.5cm}}
\toprule
Experiment & Model & Key parameters \\
\midrule
CQ-125 cross-family grounding panel & 11 LLM families & A/B all $N=3$, question-level majority vote followed by one exact two-sided McNemar test per subset; primary main-text analysis \\
GIS 125q (baseline + full) & Gemini 2.5 Flash & temp=0.0, 60s timeout \\
GIS 125q (DIN-SQL) & Gemini 2.5 Flash & temp=0.0, 4-stage pipeline \\
BIRD 108q design + 150q held-out & Gemini 2.5 Flash & temp=0.0, single-pass \\
BIRD (DIN-SQL, older 495q) & Gemini 2.5 Flash & temp=0.0, 4-stage pipeline \\
Cross-lingual 50q BIRD + 50q GIS & Gemini 2.5 Flash & LLM-translated Chinese \\
Gemini 3.5 A/B/C diagnostic & Gemini 3.5 Flash & A/B/C all N=5, question-level majority vote \\
Gemma4 fixed-host sweep & Gemma4 e2b/e4b/12b/26b/31b & Ollama 0.30.5, GGUF Q4\_K\_M, CQ-125 \\
FloodSQL-\allowbreak{}Bench smoke audit & Gemini 2.5 Flash & stratified 150q, N=1, local PostGIS port \\
\bottomrule
\end{tabular}
}
\end{table}

\paragraph{Released artefacts.} The reviewer-accessible companion repository includes the
125-question GIS benchmark with gold SQL (85 Spatial + 40 Robustness), the
108-question BIRD design set and 150-question held-out qid lists with full baseline
and full-pipeline predictions, the 100-question cross-lingual pairs (50 Chinese
BIRD + 50 Chinese GIS), the DIN-SQL runner adapted for PostgreSQL/PostGIS, the
evaluation harness with McNemar and Wilson CI tools, the balanced N=5
\texttt{gemini-3.5-flash} A/B/C summary, the six-way ablation runner,
and the SQL postprocessor and self-correction logic. Per-question predicted SQL
and execution results are included as frozen records for offline statistical
reduction; the package does not support raw-row database re-execution because the
operational rows are withheld. Revision source tables additionally include
the fixed-host Gemma4 scale sweep, the queried Ollama model metadata, and the
FloodSQL-Bench smoke-audit summary.

\paragraph{Run-time configuration.} The framework-level S1 runs use Gemini~2.5~Flash with
temperature$=0.0$, default top\_p, and default maximum output tokens. Each question
has a 60-second execution timeout; on \texttt{429 RESOURCE\_\allowbreak{}EXHAUSTED}, the runner
retries up to 3 times with 2-second initial delay. On no-SQL output, the question
is recorded with \texttt{gen\_status = no\_sql} and EX$=0$. All McNemar results
in this paper are reported as two-sided exact binomial tests on discordant pairs;
the one-sided variant is not used because we do not have a pre-registered
directional hypothesis on BIRD.

The Gemma4 sweep is documented by frozen local summaries and host metadata
rather than by a repeated-sample inferential protocol. The FloodSQL-\allowbreak{}Bench audit
uses a local PostgreSQL/\allowbreak{}PostGIS port of the benchmark and a single
\texttt{gemini-2.5-\allowbreak{}flash} stratified smoke run; it is included only to document
that the external harness was operational.

\paragraph{Reviewer access.} A reviewer-accessible companion-repository link is
provided through the editorial system for reviewer access. The primary
GIS result covers the 125-question benchmark (85 Spatial + 40 Robustness); the
BIRD evaluation is split into a 108-question design set (rule derivation) and a
150-question held-out set (independent significance).

\section*{S2. Worked example: end-to-end PostGIS grounding}

The main text \S3 (Worked example) defers the full per-stage trace to this supplement. Figure~\ref{fig:worked-example-suppl} traces a single Chongqing benchmark question through every pipeline stage. The question asks for the total length in kilometres of one-way roads; the gold SQL uses \texttt{ST\_Length} over the geography cast to return length in metres and divides by $1\,000$. Three independent grounding signals cooperate to reach the correct answer: a \emph{metric} intent, a \emph{geometry column} declaration with SRID, and an explicit \texttt{::geography} rule.

\begin{figure}[t]
\centering
\footnotesize
\begin{tabular}{@{}p{0.22\linewidth}p{0.72\linewidth}@{}}
\toprule
Stage & Content \\
\midrule
NL question ($q$) & \emph{Chinese original, transliterated:} \textit{t\v{o}ngj\`{i} su\v{o}y\v{o}u d\={a}nh\'{a}ngd\`{a}o de z\v{o}ngch\'{a}ngd\`{u} (g\={o}ngl\v{i})}. \emph{English gloss:} ``Total length in kilometres of all one-way roads.'' \\
\midrule
Intent classifier & \texttt{intent = spatial\_measurement}, \texttt{metric = length}, \texttt{unit = km}, \texttt{language = zh}. \\
\midrule
Semantic resolution (Stage~1) &
Matches \texttt{cq\_osm\_roads\_2021} as the roads table; resolves the Chinese alias for ``one-way road'' to the filter \texttt{oneway = `yes'}; identifies \texttt{geometry} column (MultiLineString, SRID $4326$, WGS~84 lat/lon). \\
\midrule
Grounded context (Stage~2) &
Schema dump for \texttt{cq\_osm\_roads\_2021}; two injected rules: (i)~\emph{metric operators on lat/lon data require \texttt{::geography} cast so that \texttt{ST\_Length} returns metres}; (ii)~\emph{unit conversion to km is a divide by $1\,000$ performed after \texttt{SUM}}. One spatial-measurement few-shot example with geography cast. \\
\midrule
Initial SQL $s^{(0)}$ &
\texttt{SELECT ROUND(SUM(ST\_Length(geometry)) / 1000.0, 2) AS length\_km}\newline
\texttt{FROM cq\_osm\_roads\_2021 WHERE oneway = 'yes';} \\
\midrule
Postprocessor + execution &
Executes, but returns a value in degree units because \texttt{::geography} was not applied; the postprocessor tags this as a ``metric unit sanity'' regression against the few-shot expectation and feeds the error message back to the grounded generator. \\
\midrule
Corrected SQL $s^{(1)}$ &
\texttt{SELECT ROUND((SUM(ST\_Length(geometry::geography)) / 1000.0)::numeric, 2)}\newline
\texttt{AS length\_km FROM cq\_osm\_roads\_2021 WHERE oneway = 'yes';} \\
\midrule
Gold SQL &
\texttt{SELECT ROUND((SUM(ST\_Length(geometry::geography)) / 1000.0)::numeric, 2)}\newline
\texttt{FROM cq\_osm\_roads\_2021 WHERE oneway = 'yes';} \\
\midrule
Outcome & Execution-based match (EX = 1); the postprocessor additionally adds the numeric cast required by PostgreSQL's \texttt{ROUND(double precision, integer)} resolution. \\
\bottomrule
\end{tabular}
\caption{Worked example: a Chongqing benchmark question traced through the full pipeline. The baseline omits the \texttt{::geography} cast; Stage~2 rule injection closes the gap.}
\label{fig:worked-example-suppl}
\end{figure}

Three implementation artefacts support this flow. The \emph{semantic layer} is a per-database YAML file declaring table aliases, geometry columns, SRIDs, geometry types, and rule identifiers. The \emph{intent rules} map classifier tags such as \texttt{spatial\_measurement} or \texttt{knn} to prompt-side rules. The \emph{postprocessor} rejects write nodes, repairs identifier case/quoting, injects \texttt{LIMIT} guards for large-table previews, and runs the retry loop above. Minimal versions of these artefacts, the few-shot top-$k$ setting, and the prompt templates are released in the reviewer-accessible companion repository.

\end{document}
